\documentclass[11pt,letterpaper]{article}
\pdfinfoomitdate=1
\pdftrailerid{}
\pdfsuppressptexinfo=15
\usepackage[margin=1in]{geometry}
\usepackage{amsmath,amssymb,amsthm}
\usepackage[T1]{fontenc}
\pdfmapfile{+lm.map}
\pdfmapfile{+cm.map}
\pdfmapfile{+symbols.map}
\usepackage{lmodern,microtype}
\usepackage[hidelinks]{hyperref}
\usepackage{xurl,enumitem}
\setlist{itemsep=3pt,topsep=5pt}
\setcounter{secnumdepth}{0}
\providecommand{\tightlist}{\setlength{\itemsep}{0pt}\setlength{\parskip}{0pt}}
\newcommand{\Log}{\operatorname{Log}}
\usepackage{booktabs,array}
% [write] notes added when the report entered the collection (batch 77)
\newenvironment{writenote}[1][]{\par\smallskip\begingroup\small
 \noindent\textbf{[write]}\if\relax\detokenize{#1}\relax\else~\textbf{#1.}\fi\ }%
 {\par\endgroup\smallskip}
\allowdisplaybreaks[2]
\hypersetup{pdftitle={Asymptotic enumeration and inversion for A122399},pdfsubject={Exact contours, all orders asymptotics, inversion, block fluctuations, and congruences}}
\title{Asymptotic enumeration and inversion for A122399}
\author{A proof and reproducible research report}
\date{1 October 2026}
\begin{document}
\maketitle
\begin{abstract}
For $a_n=\sum_k k^n k!S(n,k)$, we prove the leading equivalent already recorded in OEIS and obtain a complete fixed-order expansion with explicit coefficients. An exact vertical contour gives a branch-safe representation and a simple, nonasymptotic exponential truncation bound. The same method yields compact-direction two-size asymptotics, Gaussian block-count fluctuations, and a controlled inverse with integer-threshold envelopes. An elementary appendix proves the posted mod-prime periodicity conjecture, in the sense that $p-1$ is a period, and extends it to prime powers. The standard saddle methodology and existing leading result are explicitly credited.
\end{abstract}
\subsection{Results and attribution}\label{a122:results-and-attribution}

The leading equivalent for A122399 was posted by Vaclav Kotesovec on OEIS on August 9, 2018; its growth constant was already identified in 2013. \textbf{Neither that equivalent nor its numerical constant is claimed as a new conjecture.} This manuscript proves it and develops explicit, reproducible higher coefficients, a globally valid contour with an explicit exponentially small truncation bound, a compact-direction extension, a block-count central limit theorem, and controlled inversion. These are applications and refinements of standard analytic saddle methods, not a new general method. In particular Khera--Lundberg--Melczer (2021, preprint 2019), Section 3 and the paragraph following its equation (15), give directional asymptotics and explain higher-order expansion for the nearby but different lonesum-matrix denominator. No global novelty claim is made.

\subsection{Provenance, status and reading conventions}\label{a122:provenance}

\begin{writenote}[Added 2 October 2026, batch~77]
\emph{Provenance.} This report was built from one manuscript: batch~77,
manuscript~35 of the repository's incoming-reports channel, archive
\nolinkurl{oeis-a122399-research.zip} (main file
\nolinkurl{a122399-report.tex}, 389 lines, with an eight-page PDF that is not
shipped), arrival commit \texttt{096ee7b87}, placed in \texttt{f76fcb566}.
The text from the abstract to the end of the source list is the delivered
manuscript, unchanged except for the label prefix \texttt{a122:} and the
notes marked \textbf{[write]}. The delivered pre-layout source
\nolinkurl{proof.md} is shipped beside this article; it carries no
mathematics absent here, while Section~8 exists only in this text.
\emph{Pin.} The delivery names no ProveIt commit that it continues. Its
\nolinkurl{proof.md} (line~9) records a scoped search of the public
repository index on 1 October 2026, which found no A122399 or A317855 match
and returned other files at commit \texttt{76c0887a6}; the delivered layout
script removed that paragraph from the manuscript. The search result was
true then; the only match now is this report. The search is scoped
nonduplication evidence, as its author says, not a proof of absence.
\emph{Status.} An unrefereed research manuscript; the delivery names no
author or tool. Nothing in it is formalized in Lean or Rocq, and its
placement in the collection confers no formal status
(\hyperref[a122:elementary-congruence-appendix]{Section~7} names the one
formalized identity its proof uses). The delivered independent audit
(shipped as \nolinkurl{audit-independent_audit.md}) found no blocking
error; its script was not replayed at placement (it exceeded a 170-second
cap), while the four producer scripts were replayed with identical outputs (up to line endings).
\end{writenote}

\begin{writenote}[Reading conventions]
The manuscript reuses several letters. No symbol of the delivered text was
renamed; the table fixes the readings.

\smallskip
\noindent\begin{tabular}{@{}>{\raggedright\arraybackslash}p{0.17\linewidth}
 >{\raggedright\arraybackslash}p{0.75\linewidth}@{}}
\toprule
symbol & meanings, by section \\
\midrule
$p$ & $p(z)=1/(1+e^z)$ in Sections~2--4; $p=x/y$ in the formula for $\mu_0$
 and $p(s)$ in Section~6; a prime in Section~7. \\
$r$ & the Taylor index of $f_r$ in (3.2), (3.4); the Newton step count in
 (5.3); the prime-power exponent in (7.1); the OEIS constant
 $r=0.873702\ldots$, which is $\mu=1/y$ of (6.2), not any of the others. \\
$x,y$ & formal variables of $F(x,y;u)$ in Section~1; saddle values
 $x=X(y)$, $y=y(\lambda)$ from (3.1) on, and on the diagonal the numbers
 $x(1)$, $y(1)$ of Section~4. \\
$a$ & $a_n$ is the sequence; $a$ is a summation index in (3.4) and
 $a=y-x$ in (4.2). \\
$m$ & the red set size (with $\lambda=m/n$); $m_r$ in (3.4) are
 multiplicities. \\
$C$, $C_J$, $E_J$ & $C$ is the diagonal constant of (4.1); $C_J$ and $E_J$ in
 the proof of (5.4) are unrelated existence constants; $C(s)$ in Section~6
 has $C(0)=C$. \\
$\mu$, $\mu_0$ & mean of the block count (6.2), not a growth rate. \\
$D$ & the growth constant of (4.1): $D=1/(xy)=(1+e^{y})/y^{2}
 =3.16108865\ldots$, the constant of OEIS A317855. \\
\bottomrule
\end{tabular}

\smallskip
\noindent Section and equation numbers are the manuscript's own and are
typed by hand; equation tags such as (2.1) are fixed tags, not counters.
\end{writenote}

\subsection{1. Exact combinatorics}\label{a122:exact-combinatorics}

Put

\[
 A_{m,n}=\sum_{k=0}^m k^n k!\,S(m,k),\qquad a_n=A_{n,n},\qquad
 b_{m,n}=\frac{A_{m,n}}{m!n!}.
\]

The convention at the origin is \(0^0=1\). For positive \(m,n\), a structure consists of an ordered partition of an \(m\)-element red set into \(k\) nonempty blocks, followed by a map of an \(n\)-element blue set to these blocks. Its bivariate exponential generating function, with \(u\) marking blocks, is

\[
 F(x,y;u)=\frac1{1-u e^y(e^x-1)}.
\]

Thus the diagonal is exactly A122399, not an approximation to it. When a structure is chosen uniformly, its block count satisfies \(\Pr(K_n=k)=k^n k!S(n,k)/a_n\).

\subsection{2. Exact contour and a nonasymptotic tail bound}\label{a122:exact-contour-and-a-nonasymptotic-tail-bound}

For \(\Re z>0\), use the single analytic branch

\[
 X(z)=\Log(1+e^{-z}),\qquad p(z)=\frac1{1+e^z}=-X'(z).
\]

Here \(1+e^{-z}\) stays in the open disk of radius one about 1, and \(X(z)\ne0\). The following identities hold for every pair of integers \(m,n\ge1\) and every \(c>0\):

\[
\begin{aligned}
 b_{m,n}&=\frac1{2\pi i}\int_{c-i\infty}^{c+i\infty}
 \frac{p(z)}{X(z)^{m+1}z^{n+1}}\,dz\\
 &=\frac{n+1}{m}\frac1{2\pi i}\int_{c-i\infty}^{c+i\infty}
 X(z)^{-m}z^{-n-2}\,dz.
\end{aligned}\tag{2.1}
\]

Both integrals are absolutely convergent. Integer powers are used, so there is no unresolved power branch.

\subsubsection{Proof}\label{a122:proof}

For \(\Re h>0\) and \(n\ge1\), the absolutely and locally uniformly convergent Mittag--Leffler identity is

\[
 \frac1{n!}\sum_{k\ge1}k^n e^{-kh}
 =\sum_{j\in\mathbb Z}(h+2\pi i j)^{-n-1}.\tag{2.2}
\]

It follows, for example, by differentiating \(1/(e^h-1)\)'s symmetric partial-fraction expansion; the constant term disappears after one derivative. It also follows by Poisson summation for \(t^n e^{-ht}\mathbf1_{t\ge0}/n!\), whose Fourier coefficients are the terms on the right.

Change the \(x\)-coefficient contour to \(q=e^x-1\), with \(|q|=e^{-c}<1\); the inverse is \(x=\Log(1+q)\), analytic and one-to-one near the closed contour, and it winds once around zero. On this circle choose the continuous parametrization \(q=e^{-c+it}\), \(-\pi\le t\le\pi\), so \(h=c-it\). Insert (2.2) in the Cauchy coefficient integral. Uniform absolute convergence permits the sum--integral interchange. The shifted intervals in \(h+2\pi i j\) join to the whole vertical line; \(X\) and \(p\) are \(2\pi i\)-periodic. This gives the first identity in (2.1). Since \((X^{-m})'=m p X^{-m-1}\), integration by parts gives the second. Its boundary term tends to zero: on a fixed vertical line \(X\) is bounded away from zero, and \(z^{-n-1}\to0\). This also proves absolute convergence.

This argument does \textbf{not} treat an individual \((h(x)+2\pi i j)^{-n-1}\) as a globally analytic power series at \(x=0\). Doing so would ignore its logarithmic monodromy.

\subsubsection{Explicit truncation theorem}\label{a122:explicit-truncation-theorem}

Let \(x_c=X(c)>0\), and let \(R_T\) denote the sum of the parts of the \textbf{second} integral in (2.1) with \(|\Im z|\ge T>0\), including its prefactor. Then

\[
 |R_T|\le\frac{x_c^{-m}c^{-n}}{\pi m T}
 \left(\frac{c}{\sqrt{c^2+T^2}}\right)^n.\tag{2.3}
\]

Indeed, with \(w=X(c+it)\),

\[
 e^{-c}=|e^w-1|\le e^{|w|}-1,
 \quad\text{so}\quad |X(c+it)|\ge X(c).
\]

Consequently the modulus of the tail is bounded by \((n+1)x_c^{-m}/(\pi m)\int_T^\infty(c^2+t^2)^{-(n+2)/2}dt\). Set \(v=\sqrt{c^2+t^2}\); since \(v/\sqrt{v^2-c^2}\) decreases, that integral is at most \((c^2+T^2)^{-n/2}/((n+1)T)\), proving (2.3).

There is therefore an exact, absolutely convergent strip decomposition obtained by splitting the second contour into \((2j-1)\pi\le\Im z\le(2j+1)\pi\), \(j\in\mathbb Z\). The noncentral strips together have bound (2.3) with \(T=\pi\). They are not asserted to be canonical saddle/transseries sectors: their artificial endpoint contributions can cancel, and no individual secondary-saddle expansion is proved here. A fixed-order Poincaré expansion for the central strip has algebraic error, despite the exponentially accurate contour truncation.

\subsection{3. The unique saddle and all orders}\label{a122:the-unique-saddle-and-all-orders}

Fix \(\lambda=m/n>0\). There is a unique \(y=y(\lambda)>0\) satisfying

\[
 x=X(y)=\lambda y p(y),\qquad
 B=B(\lambda)=1+\lambda^{-1}-y(1-p(y)).\tag{3.1}
\]

Moreover \(B>1\), and \(x,y,B\) depend analytically on \(\lambda\). To prove uniqueness, \(R(y)=yp(y)/X(y)\) increases strictly from 0 to infinity. In fact \(p/X>1-p\), since \(p+(1-p)\log(1-p)>0\) for \(0<p<1\); differentiating \(\log(p/X)\) shows its derivative is positive. Thus \(R'>0\). At the root \(y(1-p)<1/\lambda\), proving \(B>1\). The implicit-function theorem supplies analyticity.

Define the analytic phase germ at \(s=0\)

\[
 f_\lambda(s)=-\log(1+s)-\lambda\log\frac{X(y(1+s))}{x}
 =\frac B2s^2+\sum_{r\ge3}f_rs^r.\tag{3.2}
\]

Here logarithms are local germs, and \(f_\lambda(0)=f'_\lambda(0)=0\). Globally along \(s=i\theta\), the exponential in the following formula means its integer-power expression, so no global choice of \(f\) is needed:

\[
 b_{m,n}=\frac{x^{-m}y^{-n}}{\lambda y}\left(1+\frac1n\right)
 \frac1{2\pi}\int_{-\infty}^{\infty}
 \frac{e^{n f_\lambda(i\theta)}}{(1+i\theta)^2}\,d\theta.\tag{3.3}
\]

The bound from Section 2 gives \(|e^{n f_\lambda(i\theta)}|\le(1+\theta^2)^{-n/2}\).

\subsubsection{Explicit coefficient formula}\label{a122:explicit-coefficient-formula}

Set \((-1)!!=1\). For \(j\ge0\), sum over nonnegative integers \(a,m_3,\ldots,m_{2j+2}\) subject to

\[
 a+\sum_{r=3}^{2j+2}(r-2)m_r=2j,
 \qquad M=a+\sum_{r=3}^{2j+2}r m_r.
\]

The integer \(M\) is necessarily even. Define

\[
 d_j=\sum
 (-1)^a(a+1)\prod_{r=3}^{2j+2}\frac{f_r^{m_r}}{m_r!}
 \frac{(-1)^{M/2}(M-1)!!}{B^{M/2}},
 \quad c_0=1,\quad c_j=d_j+d_{j-1}\ (j\ge1).\tag{3.4}
\]

The empty product gives \(d_0=1\). For every fixed integer \(J\ge0\),

\[
 \boxed{\displaystyle
 A_{m,n}=\frac{m!n!\,x^{-m}y^{-n}}{\lambda y\sqrt{2\pi n B}}
 \left(\sum_{j=0}^J\frac{c_j(\lambda)}{n^j}+O(n^{-J-1})\right).}\tag{3.5}
\]

The remainder is uniform when \(m/n\) ranges over any fixed compact subset of \((0,\infty)\). No growing-direction or boundary-uniform assertion is made.

\subsubsection{All-orders proof and remainder justification}\label{a122:all-orders-proof-and-remainder-justification}

Choose a small fixed \(\delta>0\) uniformly over the compact direction set. Outside \(|\theta|<\delta\), the modulus bound above makes the integral exponentially small; integrate the power tail exactly or use (2.3). Inside, analyticity, \(B>1\), and compactness give uniform Taylor remainder bounds and a uniform quadratic bound \(\Re f(i\theta)\le-b\theta^2\), for some \(b>0\). Set \(\theta=t/\sqrt n\). Expand the analytic amplitude \((1+s)^{-2}\) and \(\exp\{n(f(s)-Bs^2/2)\}\) through order \(n^{-J-1/2}\), keeping enough further Taylor terms to bound the remainder. On \(|t|\le n^\epsilon\), with \(\epsilon>0\) chosen sufficiently small for this fixed \(J\), Taylor's theorem bounds the remainder by \(n^{-J-1}\) times a fixed polynomial in \(|t|\) times a Gaussian with uniformly negative exponent. On \(n^\epsilon<|t|<\delta\sqrt n\), the uniform quadratic bound is superpolynomially small. Extend the polynomial--Gaussian integrals to the real line with the same accuracy. Odd terms integrate to zero. The even Gaussian moments are \((M-1)!!B^{-M/2}\), with \(i^M=(-1)^{M/2}\), giving (3.4). The prefactor \(1+1/n\) gives \(c_j=d_j+d_{j-1}\). This proves (3.5) with its uniform remainder.

This is the ordinary analytic one-dimensional saddle expansion; the exact contour is what supplies the particularly simple global tail control here.

\subsection{4. Diagonal constants and reproducible coefficients}\label{a122:diagonal-constants-and-reproducible-coefficients}

On the diagonal write \(x=x(1),y=y(1),B=2-y+x\), \(D=(xy)^{-1}\), and

\[
 C=\frac1{y\sqrt{2\pi B}}.
\]

Numerically,

\begin{itemize}
\tightlist
\item
  \(x=0.2763928914715160816083833582415360977661\)
\item
  \(y=1.144554440911905470205106411159417937429\)
\item
  \(B=1.131838450559610611403276947082118160337\)
\item
  \(D=3.161088653865428813830172202588132491726\)
\item
  \(C=0.3276282855698694814422864924105070307103\)
\end{itemize}

Thus

\[
\begin{aligned}
 a_n=\frac{C D^n(n!)^2}{\sqrt n}\Bigg[1
 &-\frac{0.0600744155748035800316371200380}{n}\\
 &+\frac{0.00759452281787645507680871266458}{n^2}\\
 &+\frac{0.00371322633571244961947936026057}{n^3}\\
 &-\frac{0.000362980958541489832391365216986}{n^4}
 +O(n^{-5})\Bigg].
\end{aligned}\tag{4.1}
\]

An exact expression for the first coefficient is

\[
 c_1=1-\frac3B+\frac{3f_4-6f_3}{B^2}-\frac{15f_3^2}{2B^3},\tag{4.2}
\]

where, setting \(a=y-x\), \(b=y-2x\),

\[
 f_3=\frac{a(b-3)}6,
\quad f_4=\frac12-\frac a2+\frac{a^2}8+\frac{ab}6
 -\frac{a(y^2-6xy+6x^2)}{24}.
\]

The numerical generator uses exact rational symbolic algebra for the phase, then 90-digit arithmetic; it does not fit these coefficients to sequence values. It prints coefficients through \(n^{-7}\), exact checks through \(n=800\), and scaled remainders. The recurrence \(T_{m,k}=k(T_{m-1,k}+T_{m-1,k-1})\), \(T_{0,0}=1\), generates \(T_{m,k}=k!S(m,k)\), so the numerical checks use independent exact integer values.

The strip-tail factor for \(T=\pi\) is \(\rho_\pi=y/\sqrt{y^2+\pi^2}<0.343\). Accordingly the exact central-contour integral differs from \(a_n/(n!)^2\) by at most \(D^n\rho_\pi^n/(\pi^2 n)\). This is a contour-evaluation statement, not an exponentially accurate bound for a fixed truncation of (4.1).

\subsection{5. Controlled inverse}\label{a122:controlled-inverse}

Define \(P_J(t)=\sum_{j=0}^J c_jt^j\). For every fixed \(J\), it is positive for all sufficiently small nonnegative \(t\). For sufficiently large real \(v\), put

\[
 H_J(v)=2\log\Gamma(v+1)+v\log D-\tfrac12\log v+\log C+\log P_J(1/v).
\]

Then

\[
 \log a_n=H_J(n)+O(n^{-J-1}),\qquad
 H'_J(v)=2\log v+\log D+O(v^{-1})>0.
\]

Consequently the large solution \(\nu_J\) of \(H_J(\nu_J)=\log a_n\) is unique and satisfies

\[
 \boxed{\nu_J=n+O(n^{-J-1}/\log n).}\tag{5.1}
\]

This follows from the mean-value theorem, using (3.5) and the derivative estimate. It is a controlled inverse along the sequence, without choosing an arbitrary exact interpolation of the integers.

For a practical Lambert-W initialization, put \(L=\log A\) and

\[
 w=W_0\left(\frac{\sqrt D}{2e}L\right),\qquad N_0=\frac{L}{2w}.
\]

If \(A=a_n\), then \(N_0=n+O(1)\), and

\[
 n=N_0-\frac{\tfrac12\log N_0+\log(2\pi C)}{2(w+1)}
 +O\left(\frac1{N_0\log N_0}\right).\tag{5.2}
\]

Indeed Stirling's expansion gives \(\log a_n=2n\log n+(\log D-2)n+\frac12\log n+\log(2\pi C)+O(1/n)\), while \(N_0\) exactly solves its first two terms. Taylor expansion around \(N_0\) gives (5.2).

\begin{writenote}[Added 2 October 2026, batch~77: an instance of repository results]
The initializer of (5.2) is an instance of the inverse-Gamma core of the
transseries-and-inversion volume
(\nolinkurl{Analysis/Transseries/docs/series-and-transseries/Transseries_And_Inversion/transseries_and_inversion.tex},
section \texttt{p6:sec:gamma}), and no novelty is claimed for the method.
The two leading Stirling terms read \(2N\log N+(\log D-2)N=L\); with
\(T=\sqrt D\,N\) and \(R=\sqrt D\,L/2\) this is the core equation
\(T(\log T-1)=R\) of \texttt{p6:eq:gamma-core-eq}, whose solution
\(w=W_0(R/e)\), \(T=R/w\) is \(w=W_0(\sqrt D L/(2e))\), \(N_0=L/(2w)\) as
above. The general power--logarithmic block is \texttt{p6:lem:core}, and the
scaled family is \texttt{p2:thm:multi-scaling}. Only the core is shared:
the lower-order terms here, \(\frac12\log N+\log(2\pi C)\), are not those of
\(\log\Gamma\), so the volume's all-orders recurrence for
\(\Gamma_+^{-1}\) does not apply as it stands; this report's own all-orders
inverse is \(\nu_J\) of (5.1).
\end{writenote}

For arbitrary algebraic inverse precision use Newton iteration

\[
 N_{r+1}=N_r-\frac{H_J(N_r)-L}{H'_J(N_r)},\tag{5.3}
\]

starting at \(N_0\). Uniformly in its \(O(1)\) starting neighborhood, \(H''_J/H'_J=O((n\log n)^{-1})\). Therefore its error relative to \(\nu_J\) after fixed \(r\) iterations is \(O((n\log n)^{-(2^r-1)})\). Taking \(2^r-1\ge J+1\) and combining with (5.1) gives the inverse error in (5.1).

For an arbitrary large target \(A\), let \(\nu=H_J^{-1}(\log A)\) and \(M(A)=\min\{n:a_n\ge A\}\). There is a constant \(E_J>0\) such that

\[
 \boxed{\left\lceil\nu-E_J\frac{\nu^{-J-1}}{\log\nu}\right\rceil
 \le M(A)\le
 \left\lceil\nu+E_J\frac{\nu^{-J-1}}{\log\nu}\right\rceil}\tag{5.4}
\]

for all sufficiently large \(A\). Indeed, the uniform sequence-value error \(|\log a_k-H_J(k)|\le C_Jk^{-J-1}\) and the lower bound on \(H'_J\) show that integers below the lower real endpoint have \(a_k<A\), and an integer at or above the upper endpoint has \(a_k\ge A\). Enlarge \(E_J\) to cover both comparisons. The sequence is strictly increasing from \(n=1\) onward: its old positive block-count summands weakly increase and additional positive structures occur.

The envelope (5.4) allows two neighboring candidates near an integer boundary, so rounding a real inverse alone is not guaranteed. Exact integer comparison settles that boundary. The constants and starting threshold in (5.4) are asymptotic existence constants, not a certified finite-input error bar unless separately bounded.

\begin{writenote}[Added 2 October 2026, batch~77]
The rounding remarks above are the general staircase facts of the same
volume, \texttt{p0:thm:staircase}: part~(1), \(M(A)=\lceil\nu\rceil\) for
the exact monotone interpolation, and part~(2), the separation condition
that fails near an integer boundary. Their generic arithmetic is formalized
in Lean as \texttt{Fabius.staircase\_ceil} and
\texttt{Fabius.staircase\_separation}
(\nolinkurl{Analysis/FabiusFunction/Lean/FabiusFunction/StaircaseInversion.lean});
those statements concern an arbitrary strictly monotone function and give
no formal status to (5.1)--(5.4), whose interpolation \(H_J\) is not exact.
\end{writenote}

\subsection{6. Block-count fluctuations}\label{a122:block-count-fluctuations}

For \(u=e^s\) near 1, replace \(X\) by \(X_s(z)=\Log(1+e^{-s-z})\), and let \(y(s)\), \(x(s)=X_s(y(s))\) solve \(x(s)=y(s)p(s)\), where \(p(s)=1/(1+e^{s+y(s)})\). Set

\[
 B(s)=2-y(s)(1-p(s)),\quad D(s)=[x(s)y(s)]^{-1},
 \quad C(s)=\frac1{y(s)\sqrt{2\pi B(s)}}.
\]

The all-orders theorem extends analytically and uniformly to a sufficiently small complex disk in \(s\). To justify this rather than assuming it, use the analytic implicit-function theorem at the nondegenerate saddle. The contour can be locally deformed through \(y(s)\) inside \(\Re(s+z)>0\). The fixed outer contour portions retain a strict exponential gap by continuity from \(s=0\), and the far tails retain uniform power bounds. The local analytic saddle expansion is therefore uniform in complex \(s\). Cauchy estimates allow differentiation of its remainder.

It follows that

\[
 \mathbb E e^{sK_n}=\frac{C(s)}{C(0)}
 \left(\frac{D(s)}{D(0)}\right)^n[1+O(n^{-1})],\tag{6.1}
\]

uniformly near zero. This is the standard analytic quasi-powers situation. A direct expansion of the logarithm of (6.1) at \(s=it/\sqrt n\), followed by the continuity theorem, gives the Gaussian limit stated below.

Let \(\mu=1/y\) and \(\sigma^2=(B-1)/(B y^2)>0\). Then

\[
 \begin{aligned}\mathbb EK_n&=\mu n+\mu_0+O(n^{-1}),\\
 \operatorname{Var}(K_n)&=\sigma^2 n+O(1),\end{aligned}
 \qquad \frac{K_n-\mu n}{\sigma\sqrt n}\Longrightarrow N(0,1),\tag{6.2}
\]

where

\[
 \mu_0=\frac{B-1}{B y}-\frac{B-1-p}{2B^2},\qquad p=x/y.
\]

For verification, envelope differentiation gives \((\log D)'=1/y\). Implicit differentiation of the saddle equation gives \(y'=1/B-1\), hence \((\log D)''=(B-1)/(By^2)\). Further, \(B'=(B-1-p)/B\) and \((\log C)'=-y'/y-B'/(2B)\), yielding \(\mu_0\).

Numerically,

\[
\begin{aligned} \mu&=0.873702433239668330496568304721,\\
 \sigma^2&=0.0889169868271170325484054553784,\\
 \mu_0&=0.144565688577761026837729354227.\end{aligned}
\]

This gives a probabilistic interpretation of the OEIS growth parametrization's constant \(r=0.873702\ldots\): it is the asymptotic fraction of occupied ordered blocks. No local limit theorem or growing-marking-range result is claimed.

\subsection{7. Elementary congruence appendix}\label{a122:elementary-congruence-appendix}

The current OEIS entry conjectures that, for a prime \(p\), the reduction of \(a_n\) for \(n\ge1\) is purely periodic with period \(p-1\). The following proves that \(p-1\) is a period; it does not assert that it is always the least period.

\textbf{Proposition.} For every prime \(p\) and integer \(r\ge1\),

\[
 a_{n+\varphi(p^r)}\equiv a_n\pmod{p^r}\qquad(n\ge r).\tag{7.1}
\]

In particular \(a_{n+p-1}\equiv a_n\pmod p\) for every \(n\ge1\).

\textbf{Proof.} Terms with \(k\ge pr\) vanish modulo \(p^r\), since \(v_p(k!)\ge\lfloor k/p\rfloor\ge r\). We may extend the sum up to \(pr-1\) even when this exceeds \(n\), since \(S(n,k)=0\) for \(k>n\). The exact inclusion-exclusion formula gives

\[
 a_n\equiv\sum_{k=0}^{pr-1}\sum_{j=0}^k
 (-1)^{k-j}\binom{k}{j}(kj)^n\pmod{p^r}.
\]

For \(n\ge r\), a term with \(p\mid kj\) is zero modulo \(p^r\). Each other base \(kj\) is a unit, so Euler's theorem leaves its power unchanged on replacing \(n\) by \(n+\varphi(p^r)\). This proves (7.1). The argument is elementary and independent of the asymptotic analysis.

\begin{writenote}[Added 2 October 2026, batch~77]
The OEIS conjecture is credited there to Peter Bala (31 May 2022), as the
delivered \nolinkurl{source_review.md} records. Neither (7.1) nor any other
statement of this report is formalized. The one nontrivial identity in its
proof, the surjection formula \(k!\,S(n,k)=\sum_{j\le k}(-1)^{k-j}\binom kj
j^n\), is formalized in Lean as
\texttt{Fabius.factorial\_mul\_stirlingSecond\_eq\_sum}
(\nolinkurl{Analysis/FabiusFunction/Lean/FabiusFunction/StirlingBasisChange.lean}),
so (7.1) is a small formalization target on top of it and Euler's theorem;
no such proof exists in the repository. The collection's report
\nolinkurl{congruences-and-valuations/secant-number-periodicity} studies the
same kind of question for the secant numbers A000364, where the analogous
OEIS conjecture (period dividing \(\varphi(m)\), pure from index one) is
false for some moduli; no result is shared, and no least period or onset is
claimed here.
\end{writenote}

\subsection{8. Questions left by this analysis}\label{a122:questions-left-by-this-analysis}

The results separate three different levels of control: an exact contour with an explicit exponential tail, a fixed-order algebraic saddle expansion, and an asymptotic inverse. Several natural extensions require additional arguments.

\begin{enumerate}
\def\labelenumi{\arabic{enumi}.}
\tightlist
\item
  \textbf{Exponentially improved expansions.} Locate and classify the complex saddles of \(z\Log(1+e^{-z})\), determine which are reached by valid contour deformations, and obtain an optimally truncated expansion with a controlled secondary-saddle remainder. The exact strip decomposition alone does not answer this question.
\item
  \textbf{Directions approaching the boundary.} Determine uniform laws when \(m/n\to0\) or \(m/n\to\infty\), including the transition scales where the compact-direction expansion ceases to be uniform. These regimes may require different saddle coordinates and block-count normalizations.
\item
  \textbf{Effective finite-input inverse bounds.} Turn the implicit all-orders constants into computable numerical bounds and starting thresholds, combining analytic Taylor estimates with the explicit contour tail. This would make the ceiling envelope a certified finite-input bracket and could reduce the need for exact comparisons away from sequence-value boundaries.
\item
  \textbf{Sharper block-count laws.} Prove a local limit theorem, explicit Edgeworth terms, or a uniform large-deviation estimate for the block count. The analytic marking germ provides the central cumulants, but the required global Fourier or tilted-contour estimates are not supplied here.
\end{enumerate}

These are limitations and possible continuations of this analysis, not claims that the corresponding general techniques or all related results are absent from the literature.

\subsection{Sources and reproducibility}\label{a122:sources-and-reproducibility}

\begin{enumerate}
\def\labelenumi{\arabic{enumi}.}
\tightlist
\item
  OEIS A122399, current entry inspected 2026-10-01: \url{https://oeis.org/A122399}
\item
  OEIS A317855: \url{https://oeis.org/A317855}
\item
  Jessica Khera, Erik Lundberg, Stephen Melczer, \emph{Asymptotic Enumeration of Lonesum Matrices}, Advances in Applied Mathematics 123 (2021), article 102118, doi:10.1016/j.aam.2020.102118; arXiv:1912.08850: \url{https://arxiv.org/html/1912.08850} (especially Section 3, Theorem 3.1, and higher-order comment after equation (15)). Its denominator is \(e^{-x}+e^{-y}-1\), different from the one here.
\item
  \nolinkurl{check_a122399.py}: exact rational phase generation, finite-sum saddle coefficients, exact-integer comparisons, and moment checks; outputs \nolinkurl{numerical_results.json}, \nolinkurl{phase_coefficients.txt}, and \nolinkurl{test_output.txt}.
\item
  \nolinkurl{export_exact_coefficients.py}: exact rational expressions through the fourth correction, including fully combined rational functions of \(x,y\) for the first two corrections; the exporter asserts that no floating-point atoms occur.
\item
  \nolinkurl{verify_contour_inverse.py}: independent direct quadrature of the finite vertical contour with the explicit tail certificate, and inverse-iteration checks.
\item
  \nolinkurl{verify_congruences.py}: finite modular checks of (7.1), including all primes below 50 and several prime powers.
\end{enumerate}

\begin{writenote}[Added 2 October 2026, batch~77]
\emph{Item~2.} The manuscript lists OEIS A317855 without saying what it
is. Checked on 2 October 2026: A317855 is the decimal expansion
\(3.161088653\ldots\) of the growth constant \(D\) of (4.1) (OEIS titles it
as a constant of the asymptotics of A122400), and the A122399 entry writes
the Kotěšovec equivalent as \(c\cdot\mathrm{A317855}^n(n!)^2/\sqrt n\) with
\(c=0.32762828\ldots=C\), its constant \(r\) being \(1/y\). The manuscript
states no other relation between A317855 and A122399.
\emph{Items~4--7 and the audit.} The scripts are shipped as
\nolinkurl{code/check_a122399.py}, \nolinkurl{code/export_exact_coefficients.py},
\nolinkurl{code/verify_contour_inverse.py} and
\nolinkurl{code/verify_congruences.py}, their recorded outputs under
\nolinkurl{data/}, and the audit as \nolinkurl{audit-independent_audit.md},
\nolinkurl{code/audit-independent_audit.py} and
\nolinkurl{data/audit-independent_audit_*}. They read and write files next
to themselves in the delivered flat layout; the report's README gives rerun
instructions that do not overwrite shipped data.
\emph{Neighbours.} Row \(n=0\) of \(A_{m,n}\) is the Fubini sequence,
\(A_{m,0}=\sum_k k!\,S(m,k)\), defined in Lean as \texttt{Fabius.fubini}
(\nolinkurl{Analysis/FabiusFunction/Lean/FabiusFunction/OrderedBell.lean});
its exact pole-lattice formula is \texttt{q2:thm:fubini} of the
transseries volume. That row lies outside the compact-direction theorem
(3.5), on the side of question~2 of Section~8, and nothing is claimed about
it here. The collection's \nolinkurl{a277364-bell-asymptotics} treats a
different Stirling sum by classical saddle methods; no theorem is shared.
\end{writenote}

The computer checks are reproducible evidence; the all-orders and error statements rely on the analytic proofs above. An independent audit and clean replay accompany this report. Conventional peer review remains appropriate before publication.

\end{document}
